\scanpage{14}

% loc: top sheet number
\mspage{2}

% loc: top
On a donc

% loc: first square
\diagram{\begin{tikzcd}[ampersand replacement=\&,column sep=large]
\underline{\operatorname{Hom}}_{\unsure{u|{}}}(I;T,T')\arrow[r]\arrow[d]\&\underline{\operatorname{Hom}}(T,T')\arrow[d]\\
{[\underline{\operatorname{Hom}}_0(I;T,T')}\arrow[r]\&{\underline{\operatorname{Hom}}_0(T,T')]}
\end{tikzcd}}

% loc: upper right of square
; avec \unsure{comm|conv}\textsuperscript{on} \unsure{en} \unsure{$\simeq$|---} si $\overline{T},\overline{T'}$ sont de \unsure{prés.} \unsure{$\varprojlim_0$}

% loc: first paragraph
% A single long diagonal stroke traverses the middle of the sheet.
Les théories \deleted{\unsure{lim}}\inserted{fini} $\varprojlim_0$ \inserted{(resp. $\varprojlim$)} forment une 2-catégorie \deleted{\unsure{les|soit}} $\unsure{\mathscr{T}|\tau|\mathcal{T}}$ (resp. $\mathscr{T}_0$), et quand $I$ est fixé, \deleted{\unsure{disons}}\inserted{via les} \deleted{\unsure{ou}} \unsure{$I$|str.}-morphismes de théories, on a \deleted{\unsure{une}} 2-catégories $\mathscr{T}^{(I)}$ ($\mathscr{T}_0^{(I)}$), et \deleted{des} \inserted{un diagramme de} foncteurs de 2-catégories

% loc: second square
\diagram{\begin{tikzcd}[ampersand replacement=\&]
\mathscr{T}\arrow[r]\arrow[d]\&\mathscr{T}^{(I)}\arrow[d,"\deleted{\unsure{\uparrow}}"']\\
\mathscr{T}_0\arrow[r]\&\mathscr{T}_0^{(I)}\end{tikzcd}}

% loc: crossed text at square
\deleted{\unsure{pour qui une fois} $\mathscr{T},\mathscr{T}$}\\
\deleted{Enfin, si on considère les théories qui sont définies par}

% loc: constructions
\uline{Constructions} \inserted{(resp. construction infinie)} pour une \deleted{\unsure{cat.}} théorie $T$ : c'est par définition un hom \deleted{dans} la théorie $i$ \deleted{$\unsure{T_0\ E_0}$} correspondant à un ens de base sans structure, donc $i_C=C$ ; donc c'est un foncteur de \unsure{2|la}-cat. 2-fibrées sur $(\mathrm{Cat}\,\varprojlim\ \text{fin})$ (resp. $(\mathrm{Cat}\,\varprojlim)$)

% loc: construction arrow
\[T\longrightarrow i\qquad\deleted{(\text{resp. }\unsure{T_0}\to i_0)}\]

% loc: category of constructions
Ces constructions \deleted{\unsure{(resp. $T$)}} forment donc une catégorie $\Gamma^T_{\deleted{\unsure{0}}}=\underline{\operatorname{Hom}}(T,i)$ (resp. \inserted{$\Gamma_0^T=$}$\underline{\operatorname{Hom}}_0(T,i)$], qui a des $\varprojlim$ finies (resp. des $\varprojlim$) et qui est munie de

% loc: evaluation
\[\begin{aligned}\Gamma^T\times T&\longrightarrow i\\(\text{resp.}\quad\Gamma_0^T\times T&\longrightarrow i_0)\end{aligned}\]

% loc: last lines
Pour $\forall$ \inserted{$C\in(\mathrm{Cat}\,\varprojlim_0)$ (resp. $\in\mathrm{Cat}\,\varprojlim$) et} $\xi\in T(C)$, le foncteur
\[\begin{aligned}u&\longmapsto u(\xi)\\{\unsure{\Pi|\Gamma}}^{T}&\longrightarrow C\quad(\text{resp. }\Gamma_0^{T}\to \unsure{C|\supset})\end{aligned}\]
commute aux $\varprojlim$ finies \deleted{\unsure{lorsque|précédentes}} (resp. quelc.)

% archive sheet number: 13
